\chapter{Herpes Simplex Virus (HSV)}

\section*{What Is This Test?}
Herpes simplex virus testing detects HSV-1 and HSV-2, double-stranded DNA viruses of the \textit{Herpesviridae} family. HSV-1 is traditionally associated with orolabial herpes ("cold sores" or "fever blisters") and is the leading cause of sporadic encephalitis in immunocompetent adults and a major cause of herpes keratitis (the leading infectious cause of corneal blindness in developed countries) \cite{whitley2001herpes}. HSV-2 is the primary cause of recurrent genital herpes. However, the epidemiologic boundaries have blurred: HSV-1 now accounts for up to 50\% of new genital herpes infections in developed countries due to changes in sexual practices (oral-genital contact). Primary infection at either site may be asymptomatic (in up to 80\% of cases for HSV-2). When symptomatic, genital herpes presents with clustered painful vesicles, pustules, or shallow ulcers on the genitals, perineum, or anus, accompanied by tender inguinal lymphadenopathy and, in primary infection, systemic symptoms (fever, malaise, myalgia). Lesions heal over 2--4 weeks. The virus establishes latency in sensory ganglia (trigeminal ganglia for HSV-1, sacral ganglia for HSV-2) and reactivates causing recurrent episodes (typically milder and shorter than primary) \cite{gupta2007genital}. HSV-2 seroprevalence in the US is approximately 12\% of persons aged 14--49; HSV-1 seroprevalence is approximately 48\%.

\section*{Alternative Names}
HSV PCR; HSV Culture; HSV Serology; HSV-1/HSV-2 Type-Specific IgG; HSV DFA; Herpes Culture; Genital Herpes Test; Tzanck Smear (historical)

\section*{Principle}
PCR: Detection of HSV-1 and HSV-2 DNA by real-time PCR from lesion swabs, CSF, or tissue. PCR is the gold standard for all sites; sensitivity >98\% for genital lesions, 96--98\% for CSF (HSV encephalitis). PCR for CSF HSV is the most common test for herpes encephalitis. Results in 2--4 hours. Viral culture: Inoculation of lesion swabs onto human diploid fibroblast cell lines (MRC-5, WI-38). Characteristic cytopathic effect (CPE) consists of cell rounding, ballooning, and formation of multinucleated syncytial giant cells, typically visible within 24--96 hours. Identification and typing (HSV-1 vs. HSV-2) by type-specific monoclonal fluorescent antibody staining. Culture sensitivity decreases significantly as lesions crust over (recurrent lesions: 30--50\% sensitivity vs. >90\% for fresh vesicles). DFA: Vesicle scraping stained with fluorescein-labeled HSV-1 or HSV-2 type-specific monoclonal antibodies. Sensitivity 80--95\% compared to PCR for fresh lesions; results within 1--2 hours. Type-specific serology: ELISA or immunoblot detecting type-specific IgG antibodies against HSV-1 glycoprotein gG-1 and HSV-2 glycoprotein gG-2 \cite{wald2002serological}. Sensitivity 95--97\%, specificity 96--98\%. IgM testing is NOT recommended for routine clinical use due to poor sensitivity, low specificity, cross-reactivity between HSV-1 and HSV-2, and inability to distinguish primary from recurrent infection. Tzanck smear (Wright or Giemsa stain): Detects multinucleated giant cells and ballooning degeneration in lesion scrapings but cannot distinguish HSV from VZV. Historical; replaced by PCR and DFA.

\section*{History}
Herpes simplex was described in ancient Greek texts (the term "herpes" means "to creep"). Transmissibility was demonstrated in the 19th century. HSV was first cultured in the 1920s by Ernest Goodpasture. The distinction between HSV-1 (above the waist) and HSV-2 (below the waist) was established in the 1960s. Acyclovir, the first effective antiherpetic agent, was developed by Gertrude Elion in the 1970s (Nobel Prize 1988) and revolutionized treatment. Type-specific serologic assays (HerpeSelect, Western blot) were developed in the 1990s. HSV PCR became widely available in the 2000s and supplanted culture and DFA as the diagnostic gold standard due to superior sensitivity and rapid turnaround. The recognition that HSV-2 seropositivity increases HIV acquisition risk 2- to 4-fold spurred interest in HSV-2 testing as an HIV prevention strategy, though HSV-2 suppressive therapy alone does not reduce HIV acquisition (demonstrated in the HPTN 039 trial). CDC does NOT recommend routine HSV-2 serologic screening in the general population.

\section*{Why This Test Is Ordered}
\begin{itemize}
  \item Diagnosis of genital herpes in patients presenting with painful genital, perineal, or perianal vesicles, pustules, or shallow ulcers; PCR of lesion specimens is the preferred diagnostic test
  \item Diagnosis of HSV encephalitis in patients presenting with acute febrile encephalopathy, altered mental status, seizures, and focal neurologic signs (particularly temporal lobe involvement with olfactory/gustatory hallucinations, personality change, or aphasia) \cite{tunkel2008management}. CSF HSV PCR is the diagnostic test of choice and should be obtained STAT
  \item Evaluation of neonatal HSV infection in infants <6 weeks with vesicular rash, seizures, sepsis syndrome, elevated liver transaminases, and/or CSF pleocytosis. HSV PCR of blood, CSF, and surface lesions/cultures (eye, mouth, rectum) is critical; disseminated neonatal HSV has 85\% mortality without treatment
  \item Diagnosis of herpes keratitis: corneal scrapings or tear film for HSV PCR or DFA in patients with dendritic corneal ulcer on slit-lamp exam
  \item Type-specific serology (HSV-1 and HSV-2 IgG) for: (a) patients with recurrent genital symptoms but repeatedly negative lesion PCR/culture; (b) determining HSV-2 status in serodiscordant couples where one partner has genital herpes, for counseling on risk reduction; (c) individuals with HIV infection (as part of the comprehensive evaluation, though routine screening remains controversial)
  \item Evaluation of esophagitis in immunocompromised patients (HSV esophagitis on upper endoscopy with PCR of biopsy tissue)
  \item Monitoring response to therapy in acyclovir-resistant HSV (foscarnet or cidofovir therapy) in immunocompromised hosts, where persistent or recurrent lesions after 7--10 days of acyclovir warrant resistance testing
\end{itemize}

\section*{Is This a Primary or Secondary Test}
HSV PCR of lesion specimens is the primary diagnostic test for mucocutaneous herpes. CSF HSV PCR is the primary test for HSV encephalitis (should be obtained STAT; acyclovir should be started before results are available). Type-specific serology is a secondary/supplementary test, NOT a primary diagnostic test for acute lesions. DFA is a primary alternative where PCR is not available.

\section*{How This Test Is Performed}
Vesicle/lesion PCR: Unroof a fresh vesicle with a sterile needle or scalpel blade, absorb the vesicular fluid with a flocked swab, and vigorously swab the base of the lesion to collect infected epithelial cells. Crusted lesions have lower sensitivity. Place swab in viral transport medium (or manufacturer-specific PCR transport medium). For CSF HSV PCR: 0.5--1 mL of CSF collected by lumbar puncture in a sterile tube; do not refrigerate. Blood: HSV PCR on EDTA whole blood or plasma for suspected neonatal HSV or disseminated HSV in immunocompromised hosts. Ocular specimens: Collected by an ophthalmologist (corneal scraping, aqueous humor). Type-specific serology: Venous blood in SST tube; serum separated within 2 hours; refrigerate or freeze for storage.

\section*{What Preparation Is Needed}
No special patient preparation. For lesion specimens, collect before application of topical antivirals, anesthetics, or drying agents. For CSF, no specific preparation beyond informed consent for lumbar puncture. For type-specific serology, no fasting required.

\section*{Contraindications and Precautions}
HSV type-specific IgG serology has several important limitations: (a) It cannot distinguish genital from orolabial infection (a positive HSV-1 IgG result does not tell you whether the patient has oral or genital HSV-1). (b) It does not indicate the duration of infection, and thus a positive result in a person with a new genital lesion could represent newly acquired or long-standing infection. (c) Seroconversion may take 3--12 weeks after infection and a negative serology during the first 3--4 weeks of a clinical episode does NOT exclude acute HSV infection. (d) The positive predictive value of HSV-2 serology in low-prevalence populations is low (approximately 50\% for an index value of 1.1--3.5), and CDC recommends confirmatory testing (HSV-2 IgG immunoblot or inhibition assay) for low-positive results. (e) IgM testing should NOT be used for diagnosis of acute HSV because it is unreliable and often falsely positive. HSV-1 seropositivity is nearly universal in adults (48--60\%) and a positive HSV-1 IgG provides essentially no diagnostic value in most clinical scenarios. CSF HSV PCR may be falsely negative in the first 24--48 hours of HSV encephalitis; if clinical suspicion is high, repeat LP with repeat PCR in 3--5 days and continue empiric acyclovir.

\section*{Risks}
Lesion sampling causes minor discomfort. Lumbar puncture carries risks of post-LP headache (10--30\%), bleeding, infection, and brain herniation (extremely rare if contraindications ruled out). There is no risk from the laboratory tests themselves. The primary clinical risk is failure to diagnose HSV encephalitis promptly: without treatment, mortality exceeds 70\%, and among survivors, severe neurologic sequelae are common. With prompt acyclovir, mortality is reduced to 10--20\%.

\section*{What Happens After the Test}
Lesion PCR results available in 2--4 hours. CSF HSV PCR available STAT (within hours). Culture: CPE may be seen in 24--48 hours; negative after 5--7 days. Serology results in 1--24 hours. For mucocutaneous HSV, first-episode treatment: acyclovir 400 mg PO TID x 7--10 days, valacyclovir 1 g PO BID x 7--10 days, or famciclovir 250 mg PO TID x 7--10 days \cite{workowski2021sexually}. Recurrent episodes (if treated): valacyclovir 500 mg PO BID x 3 days or acyclovir 800 mg PO TID x 2 days. Chronic suppressive therapy (for patients with $\geq$6 recurrences/year): valacyclovir 500--1,000 mg PO once daily, or acyclovir 400 mg PO BID. HSV encephalitis: IV acyclovir 10 mg/kg q8h for 14--21 days (do NOT use oral antivirals). Neonatal HSV: IV acyclovir 20 mg/kg q8h for 14 days (SEM disease) or 21 days (disseminated/CNS disease).

\section*{Understanding Results}

\subsection*{Below Normal}
\begin{itemize}
  \item \textbf{Lesion PCR negative for HSV-1 and HSV-2:} No HSV DNA detected. In the appropriate clinical context (vesicular lesion), this argues against HSV. However, false-negative PCR occurs with crusted lesions (<48 hours after crusting), prolonged lesion duration, or inadequate sampling. The differential includes VZV (send VZV PCR!), syphilis (chancre), chancroid, LGV, aphthous ulcers, Behçet disease, drug eruption, and contact dermatitis
  \item \textbf{CSF HSV PCR negative:} Does NOT completely exclude HSV encephalitis, especially if LP is performed in the first 24--48 hours of symptoms. If clinical suspicion for HSV encephalitis is high (fever, altered mentation, temporal lobe symptoms, MRI T2/FLAIR hyperintensity in the temporal lobes), continue acyclovir and repeat LP with repeat CSF HSV PCR in 3--5 days
  \item \textbf{Type-specific HSV-2 IgG negative:} Indicates no prior HSV-2 infection. Does NOT exclude acute HSV-2 infection if tested within the window period (seroconversion may take 3--12 weeks). A negative HSV-2 IgG in a patient with a new genital ulcer and positive lesion PCR for HSV-2 indicates primary HSV-2 infection
\end{itemize}

\subsection*{Above Normal}
\begin{itemize}
  \item \textbf{PCR positive for HSV-1:} Confirms HSV-1 infection at the sampled site. In a genital ulcer, this indicates genital HSV-1, now accounting for up to 50\% of incident genital herpes cases. Genital HSV-1 recurs far less frequently than HSV-2 (median 0--1 recurrences per year vs. 4--5 for HSV-2), and suppressive therapy is less commonly needed. CSF HSV-1 positivity in a patient with encephalitis confirms HSV-1 encephalitis; continue IV acyclovir for 21 days
  \item \textbf{PCR positive for HSV-2:} Confirms HSV-2 infection at the sampled site. Genital HSV-2 is the most common cause of recurrent genital herpes; patients should be counseled on recurrence patterns, the risk of subclinical shedding (10--20\% of days), the efficacy of suppressive therapy in reducing recurrences by 70--80\% and transmission to seronegative partners by ~50\%, and the increased risk of HIV acquisition
  \item \textbf{Type-specific HSV-2 IgG positive (index >3.5):} Confirms prior HSV-2 infection. If the patient has no history of genital herpes, they have unrecognized (asymptomatic or subclinical) infection. Counseling should address: transmission risk (subclinical shedding), availability of suppressive therapy, HIV risk, disclosure to partners, and the natural history of recurrences
  \item \textbf{Type-specific HSV-2 IgG low positive (index 1.0--3.5):} Indeterminate; may represent a false positive. The CDC recommends confirmatory testing with a second method (HSV-2 IgG inhibition assay, Western blot---available through University of Washington Clinical Virology Laboratory). True HSV-2 positivity can be verified or excluded
  \item \textbf{Tzanck smear or DFA positive:} DFA with HSV typing provides rapid diagnosis. Tzanck smear is nonspecific and should be followed by PCR or culture
\end{itemize}

\section*{Normal Results}
Lesion PCR: \textbf{HSV-1 and HSV-2 DNA not detected}. CSF PCR: \textbf{HSV DNA not detected}. Type-specific serology: \textbf{HSV-2 IgG negative (index <0.9)}. Culture: \textbf{No virus isolated at 5--7 days}. DFA: \textbf{No HSV-infected cells detected}.

\section*{Follow-up Tests}
\begin{itemize}
  \item \textbf{HSV type-specific serology (IgG):} If lesion PCR is negative in a patient with recurrent genital ulceration, serology may provide evidence of HSV-2 infection as the underlying cause
  \item \textbf{VZV PCR:} Vesicular lesions that are HSV PCR-negative should be tested for VZV (shingles/herpes zoster), which can present with similar dermatomal or disseminated vesicular eruptions
  \item \textbf{HIV testing:} All patients diagnosed with genital HSV-2 should be offered HIV testing, as HSV-2 increases HIV acquisition and transmission risk
  \item \textbf{Repeat CSF HSV PCR:} In a patient with suspected HSV encephalitis and an initial negative PCR in the first 24--48 hours, repeat LP and CSF HSV PCR at 3--5 days while continuing empiric IV acyclovir
  \item \textbf{MRI brain with contrast:} For suspected HSV encephalitis; T2/FLAIR hyperintensity in the medial temporal lobes, insula, and cingulate gyrus is characteristic; DWI may show restricted diffusion. The MRI is positive more often than CSF PCR in early disease
  \item \textbf{EEG:} In suspected HSV encephalitis, periodic lateralized epileptiform discharges (PLEDs) with temporal lobe focus are suggestive but not pathognomonic
\end{itemize}

\section*{References}
\bibliographystyle{unsrtnat}
\bibliography{references}

\clearpage
